GGSW 密文由一组 GLev 密文组成。在 GGSW 密码体制中，秘密密钥 $S$ 是由 $k$ 个多项式组成的列表，即 $S_0, S_1, ... \text{ } S_{k-1}$。GGSW 密文中的各个 GLev 密文（以 $i$ 索引）分别加密明文 $-S_0 \cdot M, -S_1 \cdot M, \ldots, -S_{k-1} \cdot M$，以及 $M$。\autoref{fig:ggsw} 直观展示了这一结构。


\subsection{加密}
\label{subsec:ggsw-enc}

\begin{tcolorbox}[title={\textbf{\tboxlabel{\ref*{subsec:ggsw-enc}} GGSW 加密}}]
$\textsf{GGSW}_{S, \sigma}^{\beta, l}(M) = \Bigl \{ \{ \textsf{GLev}_{S, \sigma}^{\beta, l}(-S_i \cdot M)  \}_{i=0}^{k-1}, \textsf{GLev}_{S, \sigma}^{\beta, l}(M) \Bigr \} \in \mathcal{R}_{\langle n, q \rangle }^{(k+1) \cdot l \cdot (k+1)}$
\end{tcolorbox}

\begin{figure}[h!]
    \centering
  \includegraphics[width=1.0\linewidth]{figures/TFHE-fig3.pdf}
  \caption{GGSW 密文结构示意图。}
  \label{fig:ggsw}
\end{figure}


\subsection{解密}

在 GLev 的消息范围与噪声条件下，解密最后一个 GLev 块即可恢复 $M$。其他块的第一行在相应误差界下解码得到 $-S_iM\bmod\beta$。要从该乘积恢复 $M\bmod\beta$，要求 $S_i$ 是 $\mathbb Z_\beta[X]/(X^n+1)$ 中的单位。非零并不充分：在 $\mathbb Z_5[X]/(X^2+1)$ 中，$X-2$ 和 $X+2$ 均非零，但 $(X-2)(X+2)=0$。使用最后一个块时，无需对 $S_i$ 求逆。



\subsection{GSW 与 RGSW}

GSW 是 $n=1$ 时的 GGSW，RGSW 是 $k=1$ 时的 GGSW。